\section{Evaluation}
\label{sec:comparative}

\textbf{Table~\ref{tab:master_bench}} compares the frozen Clay~v1.5 checkpoint with other EO and general-purpose models on 43 benchmarks. Panel~A comes from the SatML leaderboard~\cite{taylorgeospatial2026satml}, an independent harness that runs about 150 models under one protocol: linear probes on 15 classification datasets, and a UPerNet head on three segmentation datasets. We run Clay through the same harness (marked $^\dagger$). Panel~B is our own runs on benchmarks that peer models report and the leaderboard does not cover: the AlphaEarth/TESSERA suite (all 15 tasks), detection and segmentation tasks from the OlmoEarth suite~\cite{herzog2026olmoearth,olmoearth2026v13}, and sub-meter, SAR, and unseen-sensor tasks from GEO-Bench~\cite{lacoste2023geobench}, Copernicus-Bench, and TorchGeo~\cite{stewart2022torchgeo}. Each row names the best other EO model and the best general-purpose model.

\subsection{Leaderboard and Single-Date Results}
\label{sec:std_bench}

On our own single-date tasks (Table~\ref{tab:master_bench}, Panel B), Clay and the baselines we ran follow one protocol: official splits, band standardization fitted on the training split after chip-mean nodata fill, and a validation-selected $\ell_2$ logistic/Ridge probe, or a zero-decoder $1\times 1$ linear segmentation probe, over $\{\texttt{CLS},\texttt{mean}\}$. Peer cells from~\cite{olmoearth2026v13} are third-party. Classification chips under $224\,\text{px}$ are resized to $224\times 224$~\cite{corley2024resize}. Together with the leaderboard, the results fall into three regimes.

(1) \textbf{Dense tokens may help very-high-resolution segmentation (Panel B).} With a frozen $1\times 1$ linear probe and no multi-scale decoder, Clay leads \texttt{m-pv4ger-seg} and \texttt{m-chesapeake} by $4\text{--}7$ mIoU over the best peer under the same probe. Its $8\times 8$ stride gives $4\times$ as many spatial tokens as a $16\times 16$ ViT ($723$ vs. $165\,\text{GFLOPs}$~\cite{sharma2026vhrbench}). With multi-scale decoders~\cite{xiong2024dofa,waldmann2025panopticon}, \texttt{DOFA-L} / \texttt{Panopticon} score higher: $95.0/94.1$ on \texttt{m-pv4ger-seg}, $66.3/60.8$ on \texttt{m-chesapeake}, and $56.9/45.1$ on \texttt{m-cashewplant} (against $13.5/14.9$ with the linear probe).

(2) \textbf{Wavelength conditioning transfers across multispectral, SAR, and an unseen 21-band sensor (Panels A and B).} On the leaderboard's 11 Sentinel-2, Landsat, and Sentinel-1 classification datasets, Clay reaches $93\text{--}100\%$ of the best model. On four of them (\texttt{m-brick-kiln}, \texttt{forestnet}, \texttt{m-forestnet}, \texttt{treesatai}), its best entry uses RGB bands only. On 21-band Sentinel-3 OLCI (\texttt{Biomass-S3}, \texttt{LC100Cls-S3}), which it never saw in pretraining, it reaches $91\text{--}94\%$ of the best EO model. On the three segmentation datasets (Panel A), Clay with a frozen encoder reaches $92\text{--}98\%$ of the leaderboard cells, which are the best over seeds and fine-tuning modes.

(3) \textbf{RGB scene accuracy trails general backbones (Panel A).} On the four aerial RGB scene datasets (\texttt{aid}, \texttt{m-pv4ger}, \texttt{resisc45}, \texttt{ucmerced}), Clay reaches $92\text{--}98\%$ of the best model, behind general-purpose RGB backbones. A likely reason is that $0.9$ of the Stage-2 loss weight goes to reconstructing raw multi-sensor pixels. We also ran two models that are not on the leaderboard through the same harness on RGB input, and both show the same split. EmbeddingGemma~2~\cite{schechtervera2026embeddinggemma2}, a general multimodal embedder, scores $0.5\text{--}6$ points above Clay on these four datasets and $2\text{--}15$ points below it on the other 11. It ties the best general model on \texttt{aid}. FarSLIP~\cite{li2025farslip}, a ViT-B remote-sensing CLIP model, scores above Clay on these four and on \texttt{m-eurosat} and below it on the other 10. It is the best other EO model on \texttt{aid} and \texttt{resisc45}.

\subsection{Single-Date and Sensor-Agnostic}
\label{sec:aef_eval}

Table~\ref{tab:master_bench} highlights two trends in the ranking across the broader 2025--2026 evaluation suites: whether a model sees several dates, and whether it accepts any sensor. These suites are OlmoEarth-24, AlphaEarth/\allowbreak TESSERA-15, GEO-Bench-2, TerraTorch, and VHR-Bench; Table~\ref{tab:master_bench} covers part of them. In the \textbf{OlmoEarth~v1.3} and community suites, multi-temporal sequence encoders (OlmoEarth~v1.3, Galileo, TerraMind, Prithvi-EO-2.0) lead when the target depends on multi-month crop phenology (\texttt{PASTIS}, \texttt{Nandi}, \texttt{Breizh\-Crops}~\cite{olmoearth2026v13}). Dedicated optical models lead on pure-RGB texture (DINOv3 on VHR-Bench~\cite{sharma2026vhrbench}). Fixed-schema models cannot ingest sub-meter orthoimagery, $21$-band Sentinel-3, or custom band subsets without dropping channels or mapping unsupported bands into surrogate Sentinel-2 slots (\texttt{Olmo\-Earth}$^\natural$). Clay~v1.5 takes each task's own sensor. It is one of four models, with AnySat, Panopticon, and OlmoEarth, that run all $24/24$ frozen tasks in~\cite{olmoearth2026v13}. External 5-shot flood mapping across S1, S2, and PlanetScope agrees~\cite{kaushik2026flood}.

The AlphaEarth/TESSERA rows of Panel B compare single-date ViTs with precomputed $10\,\text{m}$ annual composites (TESSERA, AlphaEarth) on all 15 AlphaEarth/TESSERA tasks ($128\,\text{px}$, $5$ seeds, $100\,\text{km}$ block bootstrap). This comparison isolates two further mechanisms. First, pooling the central $2\times 2$ patch tokens ($160\times 160\,\text{m}$ around each labeled parcel) beats full-chip mean pooling at 10-shot for \texttt{Clay} on $11$ of $15$ tasks and for \texttt{OlmoEarth~v1.3} on $14$ of $15$. It is mixed for \texttt{Prithvi-EO-2.0} ($8$ of $15$; \S\ref{sec:lesson_tokens}, Lesson~2). This readout is post hoc; \texttt{CLS} was pre-specified. With it, single-date Clay has higher 10-shot point estimates than \texttt{DOFA-L}, \texttt{Prithvi-EO-2.0}, and \texttt{DINOv2-L} on $13\text{--}14$ of $15$ tasks, and the intervals exclude zero on $9\text{--}11$. Single-date \texttt{OlmoEarth~v1.3} is above Clay on $14$ of $15$ tasks ($9$ by interval).

Second, 12-month representations hold a clear advantage on seasonal crop and land-cover mapping: \texttt{OlmoEarth~v1.3}$^\diamond$, \texttt{TESSERA}, and \texttt{AlphaEarth}, plus native $64\text{-D}$ AlphaEarth on urban indicators~\cite{gong2026urban}. This is consistent with a single overpass not separating cover types whose phenology diverges across a growing season. Averaging $4$ quarterly S2 dates, with \texttt{CLS} fixed, closes a median $10\%$ of the gap to AlphaEarth at \texttt{Max} on the five tasks where the gap is at least $0.05$ ($20\%$ with concatenation). Annual composites do not apply (\texttt{n/a}) to disaster response (\texttt{Sen1Floods11}, \texttt{Vessel Detect.}). On monthly water flux (\texttt{OpenET}), the annual AlphaEarth and TESSERA samples we joined give negative $R^2$. But the published month-aligned AlphaEarth value ($0.58$) is above Clay ($0.45$).

Clay~v1.5 today appears stronger when a workflow needs one permissive Apache-2.0 checkpoint, or inputs that fixed-schema or annual systems cannot take: unseen instruments such as $21$-band Sentinel-3 OLCI, sub-meter aerial imagery with non-RGB bands, or incomplete band subsets. For multi-month phenology tasks, a multi-temporal model or annual composite is the better choice, and for RGB scene classification, a general RGB backbone such as DINOv3.
